\section{Conclusion}
On MNIST the cheapest large reduction in NFE comes from the sampler and the training-time distribution, not from the coupling. A curvature-based time grid and second-order solvers that avoid the network at $t=1$ cut the steps needed for FID 20 by 28--58\% at no training cost; the $t=1$ failure appears on every model we trained, including the nearly straight reflow models. Minibatch OT is limited by the bound $2\operatorname{tr}\Sigma^{1/2}$, and larger OT batches give no detectable gain. Logit-normal training times matched the best sampler's NFE@20. Reflow is the only change that makes one- and two-step sampling usable, because it replaces the random coupling by a deterministic one; it costs 2.7--2.8 times the training time, lowers the many-step ceiling and erases the difference between I-CFM and OT-CFM. Reflow is therefore the choice for one to four steps, a better sampler or training-time distribution from about seven steps on. Throughout, training seeds move FID far more than sampling noise, and FID and class coverage can disagree.

\paragraph{Limitations and future work.} The models are small and trained for 50 epochs without EMA, which makes the last-epoch weights noisy. Each variant has one training seed and is judged against the base models' seed spread; in particular, the logit-normal result, our largest training-side gain at many steps, rests on a single seed. FID uses ImageNet features, the DDIM reference has two seeds, and all experiments use MNIST. EMA weights, more seeds for the logit-normal model, combining it with the improved samplers, and a larger pretrained CFM model (e.g.\ on CIFAR-10) are the natural next steps.
